\section{Conclusion}
\label{sec:conclusion}

\subsection{Summary of Findings}

We have validated the Gaussian encoder used in GaussianDiffusion on MNIST across
all ten digit classes.  The key findings are:

\begin{enumerate}[leftmargin=2em]
  \item \textbf{Coordinate-inversion bug fixed.}
        The root cause of poor encoder performance was a sign error in three places
        in \texttt{src/encode.py}: the initialisation, dead-mask, and recycling
        functions all failed to account for the coordinate inversion of
        \texttt{F.affine\_grid}.  The three-line fix raises mean PSNR from
        ${\approx}26\,\text{dB}$ to $44.2 \pm 3.7\,\text{dB}$ and brings the
        alive fraction from ${\approx}60\%$ to $100\%$.

  \item \textbf{$K=70$ Gaussians is well-chosen.}
        K-sensitivity analysis confirms that $K=70$ lies in the diminishing-returns
        regime: going from $K=50$ to $K=70$ yields ${\approx}2\,\text{dB}$, but
        $K=70$ to $K=100$ yields less than $1\,\text{dB}$.

  \item \textbf{MSE loss is optimal.}
        L1\,+\,SSIM converges more slowly and achieves lower final PSNR for MNIST.
        MSE remains the recommended loss function.

  \item \textbf{Spatial structure is meaningful.}
        Gaussians cluster tightly on digit strokes, and the latent space has
        visible class-conditional structure under PCA, which is encouraging for
        the downstream generative model.

  \item \textbf{Normalization is safe.}
        All 7 parameters can be mapped to $[-1, 1]$ without clipping;
        using empirical ranges is recommended over theoretical ranges.

  \item \textbf{$\alpha$ is vestigial.}
        The opacity parameter has zero effect on rendering.
        It should be removed from the encoder output or frozen at $0.5$
        before diffusion training.

  \item \textbf{Encoding variance is tolerable.}
        Different random seeds produce parametrically different but functionally
        equivalent $\W$ tensors.  The diffusion model must handle this
        one-to-many mapping.
\end{enumerate}

\subsection{Recommendations for Diffusion Model Training}

Based on these findings we make the following recommendations before training the
DDPM GaussianTransformer:

\begin{itemize}[leftmargin=2em]
  \item \textbf{Remove or freeze $\alpha$} to avoid wasting model capacity on a
        meaningless dimension.
  \item \textbf{Use empirical per-parameter $[\min, \max]$ from the training split}
        for normalisation to maximise the effective dynamic range.
  \item \textbf{Encode the full training split at $K=70$, $T_\text{max}=3000$,
        early stop at $40\,\text{dB}$} using the corrected encoder.  All images
        tested achieve at least $37.8\,\text{dB}$; a $32\,\text{dB}$ floor can
        be set as a quality gate for the dataset.
  \item \textbf{Consider a canonical ordering} of the $K$ Gaussians
        (e.g.\ sort by position or colour) to reduce the encoding variance seen
        by the diffusion model.
  \item \textbf{Validate reconstruction on a held-out test split} before declaring
        the encoder production-ready for other datasets (Sprites, CIFAR-10).
\end{itemize}

\subsection{Limitations and Future Work}

\begin{itemize}[leftmargin=2em]
  \item The encoder is currently single-image, CPU-only, and takes ${\sim}5$--$30$\,s
        per image.  Parallelisation via SLURM job arrays (already scripted in
        \texttt{scripts/encode\_mnist.py}) is necessary for large datasets.
  \item We have only validated on MNIST.  More complex datasets (Sprites, CIFAR-10)
        will require larger $K$ and may need hyperparameter tuning.
  \item The $\alpha$ parameter should be cleanly removed from the architecture
        rather than left as dead weight.
  \item Perceptual loss (VGG features) may be beneficial for more complex images
        where MSE does not capture texture.
\end{itemize}
